% GENERATED FILE. Edit canonical lessons, then run npm run generate:books.
% canonical-source-hash: fnv1a64:5be570dbd1e90b3e
% canonical-lessons: PA-C197-routine-return, PA-C197-dinchariya, PA-C197-dinchariya-review, PA-C197-routine-label, PA-C197-payoff

\chapter{This Is a Daily Routine}
\label{ch:pa-a2-dinchariya}
\hlchaptermodality{197}

\begin{chapteropening}
\textbf{By the end of this chapter:} \emph{I can recognise and use the noun for a daily routine in one bounded identity frame.}

Recognise and produce the daily-routine noun in one bounded identity frame through listening, speaking, reading, and writing.
\end{chapteropening}

\section[\texorpdfstring{\pa{ਇਹ} · \pa{ਹੈ} (eh; hai)}{ਇਹ · ਹੈ (eh; hai)}]{Return to this and is}

\label{lesson:PA-C197-routine-return}



\begin{quote}
\emph{Hear:} \emph{eh; hai}; YOU SAY: \textbf{\pa{ਇਹ}; \pa{ਹੈ}}
\end{quote}

\begin{grammarlens}[title={Grammar Lens}]
Keep the familiar pointer first and the familiar verb last. The middle slot stays empty until the next lesson owns its new word.
\end{grammarlens}

\subsection*{Your turn}
\begin{itemize}
\raggedright
  \item \emph{Hear:} \emph{eh}; YOU SAY: \textbf{\pa{ਇਹ}}
  \item \emph{Hear:} \emph{hai}; YOU SAY: \textbf{\pa{ਹੈ}}
  \item Say the first word, pause for the empty middle, then say the last.
\end{itemize}

\begin{tcolorbox}[breakable,colback=teal!4,colframe=teal!35!black,title={Before you move on}]
Produce the two familiar words around the empty middle slot. This opener owns no new word.
\end{tcolorbox}
\section[\texorpdfstring{\pa{ਦਿਨਚਰਿਆ} (dinchariyaa)}{ਦਿਨਚਰਿਆ (dinchariyaa)}]{One noun for a daily routine}

\label{lesson:PA-C197-dinchariya}



\begin{quote}
Retrieve the familiar first and last words of the identity frame.
\end{quote}

\begin{grammarlens}[title={Grammar Lens: direct owner}]
\begin{quote}
\textbf{\pa{ਦਿਨਚਰਿਆ}} — \emph{dinchariā} — daily routine; the usual ordered activities of a day
\end{quote}

Hear and say the whole word before reading it. Keep it as one whole noun rather than inventing smaller Punjabi pieces. This lesson owns the word alone; the complete identity frame comes only after a separate retrieval.
\end{grammarlens}

\begin{tcolorbox}[breakable,skin=enhanced,colback=violet!6,colframe=violet!45!black,boxrule=0.5pt,arc=1mm,left=6pt,right=6pt,top=4pt,bottom=4pt,fonttitle=\bfseries,title={Across the family}]
An Australian Government Punjabi health page uses this exact word in an everyday routine context. A Punjabi-supported A2 lesson also uses it in its daily-routine heading. Together they support only the present-day daily-routine sense taught here; this lesson makes no claim about other specialised uses. Sources: \href{https://www.medicarementalhealthcheckin.gov.au/pa}{Australian Government Medicare Mental Health Check In}; \href{https://www.mozhilingo.com/english/punjabi/lessons/daily-routine}{MozhiLingo: Daily Routine}.

For source history, this learned Punjabi word is from Sanskrit \emph{dinacaryā}. Apte records that Sanskrit form with the sense “daily occupation, daily routine of business.” That history does not turn its parts into a productive Punjabi rule: keep the present-day Punjabi noun whole. Source: V. S. Apte, \emph{The Practical Sanskrit-English Dictionary}, p. 0813-1, \href{https://www.sanskrit-lexicon.uni-koeln.de/scans/APScan/2020/web/webtc/getword.php?key=dinacaryA\&filter=roman\&accent=no\&transLit=slp1}{Cologne Digital Sanskrit Dictionaries}.
\end{tcolorbox}

\subsection*{Your turn}
\begin{itemize}
\raggedright
  \item \emph{Hear:} \emph{dinchariā}; YOU SAY: \textbf{\pa{ਦਿਨਚਰਿਆ}}
  \item \emph{Read it:} \textbf{\pa{ਦਿਨਚਰਿਆ}}
  \item \emph{Write it:} copy \textbf{\pa{ਦਿਨਚਰਿਆ}} once, checking the whole word
\end{itemize}

\begin{tcolorbox}[breakable,colback=teal!4,colframe=teal!35!black,title={Before you move on}]
Produce the noun for “daily routine” without looking.
\end{tcolorbox}
\section[\texorpdfstring{\pa{ਦਿਨਚਰਿਆ} (dinchariyaa)}{ਦਿਨਚਰਿਆ (dinchariyaa)}]{Retrieve the new word on its own}

\label{lesson:PA-C197-dinchariya-review}



\begin{quote}
\emph{Hear:} \emph{dinchariā}; YOU SAY: \textbf{\pa{ਦਿਨਚਰਿਆ}}
\end{quote}

\subsection*{Your turn}
\begin{itemize}
\raggedright
  \item \emph{Hear:} \emph{daily routine}; YOU SAY: \textbf{\pa{ਦਿਨਚਰਿਆ}}
  \item \emph{Read it:} \textbf{\pa{ਦਿਨਚਰਿਆ}}
  \item \emph{Write it:} cover and write \textbf{\pa{ਦਿਨਚਰਿਆ}}
\end{itemize}

\begin{grammarlens}[title={Grammar Lens}]
Keep one noun and one daily-routine meaning here. Inflection, specialised uses, and a longer routine narrative stay outside this bounded lesson.
\end{grammarlens}

\begin{tcolorbox}[breakable,colback=teal!4,colframe=teal!35!black,title={Before you move on}]
Say and then write the word once from memory.
\end{tcolorbox}
\section[\texorpdfstring{\pa{ਇਹ} \pa{ਦਿਨਚਰਿਆ} \pa{ਹੈ} (eh dinchariyaa hai)}{ਇਹ ਦਿਨਚਰਿਆ ਹੈ (eh dinchariyaa hai)}]{Put the new noun into one identity frame}

\label{lesson:PA-C197-routine-label}



\begin{quote}
Retrieve \textbf{\pa{ਇਹ}}, then the new noun, then \textbf{\pa{ਹੈ}}.
\end{quote}

\begin{grammarlens}[title={Grammar Lens}]
\begin{quote}
\textbf{\pa{ਇਹ} \pa{ਦਿਨਚਰਿਆ} \pa{ਹੈ}} — \emph{eh dinchariā hai} — this is a daily routine
\end{quote}

Keep the new whole noun between familiar \textbf{\pa{ਇਹ}} and \textbf{\pa{ਹੈ}}. Familiar \textbf{\pa{ਹੈ}} stays last in the already-owned identity frame. This chapter teaches only this one label; it does not add an inflection rule or a full routine narrative.
\end{grammarlens}

\subsection*{Your turn}
\begin{itemize}
\raggedright
  \item \emph{Hear:} \emph{eh dinchariā hai}; YOU SAY: \textbf{\pa{ਇਹ} \pa{ਦਿਨਚਰਿਆ} \pa{ਹੈ}}
  \item \emph{Read it:} \textbf{\pa{ਇਹ} \pa{ਦਿਨਚਰਿਆ} \pa{ਹੈ}}
  \item \emph{Write it:} cover and write \textbf{\pa{ਇਹ} \pa{ਦਿਨਚਰਿਆ} \pa{ਹੈ}}
\end{itemize}

\begin{tcolorbox}[breakable,colback=teal!4,colframe=teal!35!black,title={Before you move on}]
Produce the complete identity frame.
\end{tcolorbox}
\section[\texorpdfstring{\pa{ਇਹ} \pa{ਦਿਨਚਰਿਆ} \pa{ਹੈ} (eh dinchariyaa hai)}{ਇਹ ਦਿਨਚਰਿਆ ਹੈ (eh dinchariyaa hai)}]{A bounded daily-routine noun payoff}

\label{lesson:PA-C197-payoff}



\begin{quote}
Retrieve the new noun alone before rebuilding its identity frame.
\end{quote}

\subsection*{Your turn: four skills}
\begin{itemize}
\raggedright
  \item \emph{Hear:} \emph{eh dinchariā hai}; YOU SAY: \textbf{\pa{ਇਹ} \pa{ਦਿਨਚਰਿਆ} \pa{ਹੈ}}
  \item \emph{Read it:} \textbf{\pa{ਇਹ} \pa{ਦਿਨਚਰਿਆ} \pa{ਹੈ}}
  \item \emph{Write it:} cover the line and write \textbf{\pa{ਇਹ} \pa{ਦਿਨਚਰਿਆ} \pa{ਹੈ}}
\end{itemize}

\begin{tcolorbox}[breakable,colback=teal!4,colframe=teal!35!black,title={Before you move on}]
Say the new noun, then produce the complete identity frame.

This is the twentieth prerequisite tranche exposed by Punjabi A2 mock drafting. It does not unblock Punjabi A2 mock 1 by itself. It is not an A2 mock, not A2 readiness or attainment, and not external or human validation.
\end{tcolorbox}
